\Needspace{8\baselineskip}
\Article{185}
Die Anzahl aller reduzierten Formen mit gegebener Determinante $D$ ist stets endlich; die Formen selbst lassen sich auf zwei Arten finden. Bezeichnen wir alle reduzierten Formen mit Determinante $D$ allgemein durch $(a,b,c)$, sodass sämtliche Werte von $a,b,c$ zu bestimmen sind.

\textit{Erste Methode.} Für $a$ nehme man alle Zahlen (sowohl positiv als auch negativ), die kleiner als $2\surd D$ sind und für die $D$ ein quadratischer Rest ist. Für jedes $a$ setze man $b$ gleich allen positiven Werten des Ausdrucks $\surd D\ (\mathrm{mod.}\ a)$, die zwischen $\surd D$ und $\surd D\mp a$ liegen; $c$ setze man dagegen für jedes bestimmte Wertepaar von $a,b$ $=\frac{bb-D}{a}$. Ergeben sich dabei Formen, in denen $\pm a$ außerhalb des Intervalls zwischen $\surd D+b$ und $\surd D-b$ liegt, so sind sie zu verwerfen.

\textit{Zweite Methode.} Für $b$ nehme man alle positiven Zahlen kleiner als $\surd D$. Für jedes $b$ zerlege man $bb-D$ auf alle möglichen Arten in zwei Faktoren, die ohne Berücksichtigung ihres Vorzeichens zwischen $\surd D+b$ und $\surd D-b$ liegen, und setze den einen $=a$, den anderen $=c$. Jede Zerlegung liefert offenbar zwei Formen, da jeder Faktor sowohl $=a$ als auch $=c$ gesetzt werden muss.
\par
\SourceLeaf{172}
\textit{Beispiel.} Sei $D=79$; dann sind die zweiundzwanzig Werte von $a$: ${\mp1}$, $2$, $3$, $5$, $6$, $7$, $9$, $10$, $13$, $14$, $15$. Daraus ergeben sich neunzehn Formen:
\[
\begin{gathered}
(1,8,-15),\quad (2,7,-15),\quad {(3,8,-5)},\quad (3,7,-10),\\
(5,8,-3),\quad (5,7,-6),\quad (6,7,-5),\quad (6,5,-9),\\
(7,4,-9),\quad (7,3,-10),\quad (9,5,-6),\quad (9,4,-7),\\
(10,7,-3),\quad (10,3,-7),\quad (13,1,-6),\quad (14,3,-5),\\
(15,8,-1),\quad (15,7,-2),\quad (15,2,-5)
\end{gathered}
\]
und ebenso viele weitere, die durch Umkehrung der Vorzeichen der äußeren Glieder entstehen, nämlich $(-1,8,15)$, $(-2,7,15)$ usw., insgesamt also achtunddreißig. Davon sind aber sechs zu verwerfen: $(\pm13,1,\mp6)$, $(\pm14,3,\mp5)$, $(\pm15,2,\mp5)$; die übrigen zweiunddreißig umfassen alle reduzierten Formen. Die zweite Methode liefert dieselben Formen in folgender Reihenfolge\HistoricalNote{*)}{Für $b=1$ lässt sich ${-78}$ nicht in zwei Faktoren zerlegen, die ohne Berücksichtigung ihres Vorzeichens zwischen $\surd79+1$ und $\surd79-1$ liegen. Daher ist dieser Wert zu überspringen, aus demselben Grund auch die Werte $2$ und $6$.}:
\[
\begin{gathered}
(\pm7,3,\mp10),\quad (\pm10,3,\mp7),\quad {\pm7,4,\mp9)},\quad (\pm9,4,\mp7),\\
(\pm6,5,\mp9),\quad (\pm9,5,\mp6),\quad (\pm2,7,\mp15),\quad (\pm3,7,\mp10),\\
(\pm5,7,\mp6),\quad (\pm6,7,\mp5),\quad (\pm10,7,\mp3),\quad (\pm15,7,\mp2),\\
(\pm1,8,\mp15),\quad (\pm3,8,\mp5),\quad (\pm5,8,\mp3),\quad (\pm15,8,\mp1).
\end{gathered}
\]

\Needspace{8\baselineskip}
\Article{186}
Sei $F$ eine reduzierte Form mit Determinante $D$ und $F'$ ihre nach rechts benachbarte reduzierte Form; dieser sei wiederum die reduzierte Form $F''$ nach rechts benachbart, die reduzierte Form $F'''$ der Form $F''$ usw. Dann sind alle Formen $F',F'',F'''$ usw. vollständig bestimmt und sowohl untereinander als auch zu $F$ eigentlich äquivalent. Da aber die Anzahl aller reduzierten Formen mit gegebener Determinante endlich ist, können nicht alle Formen der unendlichen Folge $F,F',F''$ usw. verschieden sein. Nehmen wir an, $F^m$ und $F^{m+n}$ seien identisch. Dann sind $F^{m-1}$, $F^{m+n-1}$ reduzierte Formen, die derselben reduzierten Form nach links benachbart sind, und daher identisch; ebenso sind $F^{m-2}$ und $F^{m+n-2}$ usw. und schließlich $F$ und $F^n$ identisch. Wird also die Folge $F,F',F''$ usw. weit genug fortgesetzt, so muss schließlich die erste Form $F$ wiederkehren. Nehmen wir an, $F^n$ sei die erste mit $F$ identische Form, seien also alle $F',F''\mathinner{\ldotp\ldotp\ldotp\ldotp} F^{n-1}$ von $F$ verschieden, so sieht man leicht, dass sämtliche Formen $F,F',F''\mathinner{\ldotp\ldotp\ldotp\ldotp} F^{n-1}$ verschieden sind. Diese Gesamtheit von \SourceLeaf{173}Formen nennen wir die \textit{Periode der Form} $F$. Setzt man die Folge über die letzte Form der Periode hinaus fort, so erscheinen dieselben Formen $F,F',F''$ usw. erneut, und die ganze unendliche Folge $F,F',F''$ usw. besteht aus dieser unendlich oft wiederholten Periode von $F$.

Die Folge $F,F',F''$ usw. lässt sich auch rückwärts fortsetzen, indem man der Form $F$ die ihr nach links benachbarte reduzierte Form ${}'F$, dieser wiederum die ihr nach links benachbarte reduzierte Form ${}''F$ usw. voranstellt. So entsteht eine nach \textit{beiden Seiten} unendliche Folge von Formen:
\[
\mathinner{\ldotp\ldotp\ldotp\ldotp} {}'''F,\ {}''F,\ {}'F,\ F,\ F',\ F'',\ F'''\mathinner{\ldotp\ldotp\ldotp\ldotp}
\]
Man sieht leicht, dass ${}'F$ mit $F^{n-1}$, ${}''F$ mit $F^{n-2}$ usw. identisch ist; damit besteht die Folge auch nach links aus der unendlich oft wiederholten Periode von $F$.

Ordnet man den Formen $F,F',F''$ usw., ${}'F$, ${}''F$ usw. die Indizes $0$, $1$, $2$ usw., $-1$, $-2$ usw. zu, allgemein der Form $F^m$ den Index $m$ und der Form ${}^mF$ den Index $-m$, so sind offenbar \textit{je zwei Formen der Folge identisch oder verschieden, je nachdem ihre Indizes modulo $n$ kongruent oder inkongruent sind.}

\textit{Beispiel.} Als Periode der Form $(3,8,-5)$, deren Determinante $=79$ ist, findet man ${3,8,-5)}$, $(-5,7,6)$, $(6,5,-9)$, $(-9,4,7)$, $(7,3,-10)$, $(-10,7,3)$. Nach der letzten erscheint erneut $(3,8,-5)$. Hier ist also $n=6$.

\Needspace{8\baselineskip}
\Article{187}
Hier folgen einige allgemeine Bemerkungen über diese Perioden.

1) Schreibt man die Formen $F,F',F''$ usw.; ${}'F,{}''F,{}'''F$ usw. als
\[
\begin{gathered}
{a,b,-a')},\quad (-a',b',a''),\quad (a'',b'',-a''')\text{ usw.};\\
(-{}'a,{}'b,a),\quad ({}''a,{}''b,-{}'a),\quad (-{}'''a,{}'''b,{}''a)
\end{gathered}
\]
so haben alle $a,a',a'',a'''$ usw., ${}'a,{}''a,{}'''a$ usw. \textit{dasselbe Vorzeichen} (Artikel~184,~1), während alle $b,b',b''$ usw., ${}'b,{}''b$ usw. positiv sind.

\Needspace{7\baselineskip}
2) Daraus folgt, dass die Zahl $n$ (die Anzahl der Formen, aus denen die Periode von $F$ besteht) stets \textit{gerade} ist. Denn das erste Glied jeder Form $F^m$ dieser Periode hat offenbar dasselbe Vorzeichen wie das erste Glied $a$ von $F$, wenn $m$ gerade ist, und das entgegengesetzte, wenn $m$ ungerade ist. Da $F^n$ und $F$ identisch sind, muss $n$ gerade sein.
\SourceLeaf{174}

\Needspace{14\baselineskip}
3) Nach Artikel~184,~6 lautet der Algorithmus zur Bestimmung der Zahlen $b',b'',b'''$ usw., $a'',a'''$ usw.:
\[
\begin{array}{lll}
b'\equiv-b\ (\mathrm{mod.}\ a')&\text{zwischen }\surd D\text{ und }\surd D\mp a';&a''=\dfrac{D-b'b'}{a'}\\[6pt]
b''\equiv-b'\ (\mathrm{mod.}\ a'')&\mathinner{\ldotp\ldotp\ldotp\ldotp\ldotp\ldotp\ldotp\ldotp\ldotp}\quad\surd D\mp a'';&a'''=\dfrac{D-b''b''}{a''}\\[6pt]
b'''\equiv-b''\ (\mathrm{mod.}\ a''')&\mathinner{\ldotp\ldotp\ldotp\ldotp\ldotp\ldotp\ldotp\ldotp\ldotp}\quad\surd D\mp a''';&a''''=\dfrac{D-b'''b'''}{a'''}
\end{array}
\]
\centerline{usw.}
In der zweiten Spalte sind die oberen oder unteren Vorzeichen zu nehmen, je nachdem $a,a',a''$ usw. positiv oder negativ sind. Anstelle der Formeln in der dritten Spalte kann man auch die folgenden verwenden, die bequemer sind, wenn $D$ groß ist:
\[
\begin{aligned}
a''&=\frac{b+b'}{a'}(b-b')+a\\
a'''&=\frac{b'+b''}{a''}(b'-b'')+a'\\
a''''&=\frac{b''+b'''}{a'''}({b''}-b''')+a''
\end{aligned}
\]
\centerline{usw.}

4) Jede Form $F^m$ in der Periode von $F$ hat im Wesentlichen dieselbe Periode wie $F$. Diese lautet nämlich $F^m$, $F^{m+1}\mathinner{\ldotp\ldotp\ldotp\ldotp} F^{n-1}$, $F,F'\mathinner{\ldotp\ldotp\ldotp\ldotp} F^{m-1}$; sie enthält dieselben Formen in derselben Reihenfolge wie die Periode von $F$ und unterscheidet sich von ihr nur durch Anfang und Ende.

\Needspace{30\baselineskip}
5) Daher lassen sich alle reduzierten Formen mit derselben Determinante $D$ in Perioden \textit{aufteilen}. Man nehme eine beliebige dieser Formen, $F$, und bestimme ihre Periode $F,F',F''\mathinner{\ldotp\ldotp\ldotp\ldotp} F^{n-1}$, die wir mit $P$ bezeichnen. Umfasst sie noch nicht alle reduzierten Formen mit Determinante $D$, so sei $G$ eine darin nicht enthaltene Form und $Q$ ihre Periode. Dann können $P$ und $Q$ offenbar keine Form gemeinsam haben; andernfalls müsste auch $G$ in $P$ enthalten sein, und die Perioden würden vollständig zusammenfallen. Sind mit $P$ und $Q$ noch nicht alle reduzierten Formen erfasst, so liefert eine der fehlenden Formen, $H$, eine dritte Periode $R$, die weder mit $P$ noch mit $Q$ eine Form gemeinsam hat. So können wir fortfahren, bis alle reduzierten Formen erfasst sind. \textit{Zum Beispiel} teilen sich alle reduzierten Formen mit Determinante $79$ in sechs Perioden auf:
\SourceLeaf{175}
\[
\begin{array}{rl}
\mathrm{I.}&(1,8,-15),\quad (-15,7,2),\quad (2,7,-15),\quad (-15,8,1).\\
\mathrm{II.}&(-1,8,15),\quad (15,7,-2),\quad (-2,7,15),\quad (15,8,-1).\\
\mathrm{III.}&(3,8,-5),\quad (-5,7,6),\quad (6,5,-9),\\
&(-9,4,7),\quad (7,3,-10),\quad (-10,7,3).\\
\mathrm{IV.}&(-3,8,5),\quad (5,7,-6),\quad (-6,5,9),\\
&(9,4,-7),\quad (-7,3,10),\quad (10,7,-3).\\
\mathrm{V.}&(5,8,-3),\quad (-3,7,10),\quad (10,3,-7),\\
&(-7,4,9),\quad (9,5,-6),\quad (-6,7,5).\\
\mathrm{VI.}&(-5,8,3),\quad (3,7,-10),\quad (-10,3,7),\\
&(7,4,-9),\quad (-9,5,6),\quad (6,7,-5).
\end{array}
\]

6) Formen, die aus denselben Gliedern in umgekehrter Reihenfolge bestehen, etwa $(a,b,-a')$, $(-a',b,a)$, nennen wir \textit{assoziierte Formen}. Aus Artikel~184,~7 sieht man leicht: Hat die reduzierte Form $F$ die Periode $F,F',F''\mathinner{\ldotp\ldotp\ldotp\ldotp} F^{n-1}$, ist der Form $F$ die Form $f$ assoziiert und sind den Formen $F^{n-1},F^{n-2}\mathinner{\ldotp\ldotp\ldotp\ldotp} F'',F'$ jeweils die Formen $f',f''\mathinner{\ldotp\ldotp\ldotp\ldotp} f^{n-2},f^{n-1}$ assoziiert, so ist die Periode von $f$ gleich $f,f',f''\mathinner{\ldotp\ldotp\ldotp\ldotp} f^{n-2},f^{n-1}$. Sie besteht also aus ebenso vielen Formen wie die Periode von $F$. Die Perioden assoziierter Formen nennen wir \textit{assoziierte Perioden}. In unserem Beispiel sind die Perioden III und VI sowie IV und V assoziiert.

\Needspace{24\baselineskip}
7) Es kann aber auch die Form $f$ selbst in der Periode der ihr assoziierten Form $F$ vorkommen, wie in unserem Beispiel bei den Perioden I und II. Dann stimmt die Periode von $F$ mit der von $f$ überein, oder \textit{die Periode von $F$ ist sich selbst assoziiert}. In diesem Fall findet man in dieser Periode zwei ambige Formen. Nehmen wir nämlich an, die Periode von $F$ bestehe aus $2n$ Formen, seien also $F$ und $F^{2n}$ identisch; ferner sei $2m+1$ der Index von $f$ in der Periode von $F$\HistoricalNote{*)}{Dieser Index muss ungerade sein, da die ersten Glieder der Formen ${F},f$ offenbar entgegengesetzte Vorzeichen haben (\textit{siehe oben}, $2$).}, seien also $F^{2m+1}$ und $F$ assoziiert. Dann sind auch $F'$ und $F^{2m}$, ferner $F''$ und $F^{2m-1}$ usw. und damit auch $F^m$ und $F^{m+1}$ assoziiert. Sei $F^m=(a^m,b^m,-a^{m+1})$, $F^{m+1}=(-a^{m+1},b^{m+1},a^{m+2})$. Dann gilt $b^m+b^{m+1}\equiv0\ (\mathrm{mod.}\ a^{m+1})$; nach der Definition assoziierter Formen gilt aber $b^m=b^{m+1}$ und somit $2b^{m+1}\equiv0\ (\mathrm{mod}\ a^{m+1})$, die Form $F^{m+1}$ ist also ambig. --- Ebenso sind $F^{2m+1}$ und $F^{2n}$ assoziiert; daher auch $F^{2m+2}$ und $F^{2n-1}$, $F^{2m+3}$ und $F^{2n-2}$ usw. und schließlich $F^{m+n}$ und $F^{m+n+1}$. Die letztere ist ambig, wie ein ähnlicher Schluss leicht zeigt. Da aber $m+1$ und $m+n+1$ modulo $2n$ inkongruent sind, sind die Formen $F^{m+1}$ und $F^{m+n+1}$ nicht identisch (Artikel~186, wo $n$ dasselbe bezeichnet wie hier $2n$). So sind in I die Formen $(1,8,-15)$, $(2,7,-15)$ und in II die Formen $(-1,8,15)$, $(-2,7,15)$ ambig.
\SourceLeaf{176}

8) Umgekehrt ist \textit{jede Periode, in der eine ambige Form vorkommt, sich selbst assoziiert.} Denn ist $F^m$ eine reduzierte ambige Form, so sieht man leicht, dass die ihr assoziierte Form (die ebenfalls reduziert ist) zugleich ihre nach links benachbarte Form ist, \textit{das heißt}, $F^{m-1}$ und $F^m$ sind assoziiert. Dann ist aber die ganze Periode sich selbst assoziiert. --- Daraus folgt, dass \textit{eine Periode nicht genau eine ambige Form enthalten kann.}

9) Sie kann aber auch \textit{nicht mehr als zwei enthalten.} Nehmen wir an, die Periode von $F$, die aus $2n$ Formen besteht, enthalte drei ambige Formen $F^\lambda,F^\mu,F^\nu$ mit den jeweiligen Indizes $\lambda,\mu,\nu$, sodass $\lambda,\mu,\nu$ verschiedene Zahlen zwischen $0$ und $2n-1$ einschließlich der Grenzen sind. Dann sind die Formen $F^{\lambda-1}$ und $F^\lambda$, ebenso $F^{\lambda-2}$ und $F^{\lambda+1}$ usw. und schließlich $F$ und $F^{2\lambda-1}$ assoziiert. Aus demselben Grund sind $F$ und $F^{2\mu-1}$ sowie $F$ und $F^{2\nu-1}$ assoziiert. Daher sind $F^{2\lambda-1}$, $F^{2\mu-1}$, $F^{2\nu-1}$ identisch und die Indizes $2\lambda-1$, $2\mu-1$, $2\nu-1$ modulo $2n$ kongruent, folglich auch $\lambda\equiv\mu\equiv\nu\ (\mathrm{mod.}\ n)$. \textit{Das ist widersinnig}, da zwischen $0$ und $2n-1$ offenbar keine drei verschiedenen Zahlen liegen können, die modulo $n$ kongruent sind.
